% !TeX program = lualatex
% Plano de aula: versão para a coordenação, os docentes e o arquivo. Esqueleto comum com o deck: aula/03-slides/argumento.md.
% Compila com pdfLaTeX ou LuaLaTeX. Troque [final] por nada enquanto houver pendências (\pendente, \averificar, \figura
% aparecem em vermelho no rascunho e viram erro no modo final).
\documentclass[11pt]{article}
\usepackage{aula-geo}
\graphicspath{{../../../fontes/provas-recortes/}}
../estilo/metadados.tex
\definir\tipodocumento{Plano de aula}
\definir\versao{v0}

\begin{document}
\blocotitulo

\section*{Identificação}
\begin{tabularx}{\linewidth}{@{}>{\sffamily\small\color{fgsoft}}l X@{}}
Curso & Pro Enem (pré-vestibular para o ENEM) \\
Área e disciplina & \pendente{área; disciplina} \\
Turma & \pendente{turma} \\
Data & \auladata \\
Duração & 50 minutos \\
Docentes & \aulaprofessor \\
Tema & \pendente{tema da ementa} \\
Pergunta norteadora & \pendente{pergunta-guia} \\
Modalidade & Aula expositiva dialogada, com questões do ENEM resolvidas com a turma \\
\end{tabularx}

\section{Apresentação e justificativa}
\pendente{por que esta aula cabe no ENEM e o que o diagnóstico da turma pede}

\section{Objetivos}
\begin{itemize}
\item \pendente{objetivo 1}
\end{itemize}

\section{Roteiro}
\begin{roteiro}
0--5 & Abertura & \pendente{pergunta-guia e gancho} & slide \texttt{titulo} \\
5--40 & Desenvolvimento & \pendente{blocos do argumento} & slides \\
40--50 & Questões do ENEM & \pendente{4 questões conferidas} & recortes \\
\end{roteiro}

\section{Questões}
% \questaoenem{2013_CH_001}{2013}{1}{E}   % id do recorte em fontes/provas-recortes/, ano, nº oficial, gabarito
\begin{questaoautoral}{A}
\pendente{enunciado autoral}
\alternativas
\item alternativa A
\item alternativa B
\item alternativa C
\item alternativa D
\item alternativa E
\fimalternativas
\end{questaoautoral}
\imprimirgabarito

\section{Referências}
\pendente{só o que está em fontes/checagem.md com status ok}
\end{document}
